% Written By Enoch Yu
% Downloaded from archive.enochyu.com

\documentclass[12pt]{report}
\usepackage{graphicx}
\usepackage{amsmath,amsthm,amssymb}
\usepackage[font=small,labelfont=bf]{caption}
\usepackage{tikz}
\usepackage{float}
\usepackage[margin=1in]{geometry}
\usepackage{gensymb}
\usepackage{hyperref}
\hypersetup{
    colorlinks=true,
    linkcolor=blue,
    filecolor=magenta,      
    urlcolor=cyan,
    pdftitle={Ancient File 3},
    pdfpagemode=FullScreen,
}
\usepackage{fancyhdr}
\pagestyle{fancy}
\fancyhead[R]{Enoch Yu}
\newtheorem{theorem}{Theorem}[section]
\newtheorem{lemma}[theorem]{Lemma}
\newtheorem{proposition}{Proposition}
\newtheorem{corollary}{Corollary}[theorem]

\title{Ancient File 3}
\author{Enoch Yu}
\date{March 2025}

\begin{document}

\maketitle

\begin{center}
\small\em
Hello! My math group and I made bunch of fun problems
in the past, and I wrote solutions for some problems.
Apparently, I lost the problem sets and only have the
solutions :( Here's what I have!
\end{center}

\chapter*{Problem 1}
\textbf{Key Word}: $\frac{1}{n(n+k)}=\frac{1}{k}(\frac{1}{n}-\frac{1}{n+k})$
\\\\
\setlength{\jot}{8pt}
\begin{align*}
&\frac{1}{3}-\frac{1}{4}+\frac{1}{15}-\frac{1}{12}+\dots+\frac{1}{1155}-\frac{1}{612}+\frac{1}{1295}-\frac{1}{684} \\
=& \frac{1}{3}+\frac{1}{15}+\dots+\frac{1}{1155}+\frac{1}{1295}-\frac{1}{4}-\frac{1}{12}\dots-\frac{1}{612}-\frac{1}{684} \\
=& \frac{1}{1\cdot3}+\frac{1}{3\cdot5}+\dots+\frac{1}{33\cdot35}+\frac{1}{35\cdot37}-\frac{2}{2\cdot4}-\frac{2}{4\cdot6}\dots-\frac{2}{34\cdot36}-\frac{2}{36\cdot38} \\
=& \frac{1}{2}(\frac{1}{1}-\frac{1}{3}+\frac{1}{3}-\frac{1}{5}+\dots+\frac{1}{33}-\frac{1}{35}+\frac{1}{35}-\frac{1}{37})-2[\frac{1}{2}(\frac{1}{2}-\frac{1}{4}+\frac{1}{4}-\frac{1}{6}+\dots+\frac{1}{34}-\frac{1}{36}+\frac{1}{36}-\frac{1}{38})] \\
=& \frac{1}{2}(\frac{1}{1}-\frac{1}{37})-2[\frac{1}{2}(\frac{1}{2}-\frac{1}{38})] \\
=& \frac{1}{2}(\frac{1}{1}-\frac{1}{37})-(\frac{1}{2}-\frac{1}{38}) \\
=& \frac{1}{2}\cdot\frac{36}{37}-(\frac{1}{2}-\frac{1}{38}) \\
=& \frac{18}{37}-\frac{18}{38} \\
=& \frac{18(38-37)}{37\cdot38} \\
=& \frac{18}{37\cdot38} \\
=& \frac{9}{703}
\end{align*}
\\\\
\\\\
Conclusion: \textbf{The computed value is $\frac{9}{703}$}.


\chapter*{Problem 5}
\textbf{Key Word}: Similar Triangle
\\\\
First and foremost, the given ratios of diverse length may be altered with least common multiple to avoid fraction.
\\\\
$LCM(4+3, 3+1,5+3)=LCM(7,4,8)=56$
\\\\
Thereby, $\overline{AD}:\overline{DC}=4:3=32:24, \overline{CF}:\overline{FB}=3:1=42:14 \text{ and }\overline{BE}:\overline{EA}=5:3=35:21$. \\
\begin{figure}[H]
    \includegraphics[scale=0.3]{Images/DIAGRAM1.png}
    \caption{The circles, triangles and rectangles indicate that the numbers enclosed are ratios.}
\end{figure}
To exploit the property of similar triangles, the three lines parallel to the sides (colored in red) may be drawn.
\begin{figure}[H]
    \includegraphics[scale=0.3]{Images/DIAGRAM2.png}
\end{figure}
Using the property of similar triangles, more specified ratios of the lengths are feasibly found.
\begin{align*}
    &\overline{AE}:\overline{EB}=\overline{AG}:\overline{GC}=21:35 \text{ } (\because \triangle AGE \sim \triangle ACB) \\
    &\overline{AG}:\overline{GD}:\overline{DC}=21:11:24
    \\\\
    &\overline{CD}:\overline{DA}=\overline{CI}:\overline{IB}=24:35 \text{ } (\because \triangle CDI \sim \triangle CAB) \\
    &\overline{CI}:\overline{IF}:\overline{FB}=24:18:14
    \\\\
    &\overline{BF}:\overline{FC}=\overline{BH}:\overline{HA}=14:42 \text{ } (\because \triangle BFH \sim \triangle BCA) \\
    &\overline{BH}:\overline{HE}:\overline{EA}=14:21:21
\end{align*}
If $x$ is the area of $\triangle ABC$, the following diagram manifests the specific areas within the triangle.
\begin{figure}[H]
    \includegraphics[scale=0.3]{Images/DIAGRAM3.png}
    \caption{Using the properties of similar triangle, the specific areas within $\triangle ABC$ may be found.}
\end{figure}
\begin{align*}
    &\frac{9}{64}x+\frac{33}{448}x:\frac{9}{49}x+\frac{27}{196}x:\frac{3}{32}x+\frac{1}{16}x \\
    =& \frac{63}{448}x+\frac{33}{448}x:\frac{36}{196}x+\frac{27}{196}x:\frac{3}{32}x+\frac{2}{32}x \\
    =& \frac{96}{448}:\frac{63}{196}:\frac{5}{32} \\
    =& \frac{3}{14}:\frac{9}{28}:\frac{5}{32} \\
    =& \frac{48}{14\cdot16}:\frac{72}{28\cdot8}:\frac{35}{32\cdot7} \\
    =& 48:72:35
\end{align*}
\\\\
\\\\
Conclusion: \textbf{$48:72:35$ is the ratio of the areas of $\triangle ADE, \triangle DCF, \text{and } \triangle BEF$}.


\chapter*{Problem 7}

\noindent
Key Word: $\frac{a}{b}=\frac{c}{d}=\frac{a+c}{b+d}$
\\\\
\begin{align*}
    k&=\frac{ab}{2bc-3a+1}=\frac{bc}{2cd-3b-1}=\frac{cd}{2da-3c+1}=\frac{da}{2ab-3d-1} \\
    &=\frac{(ab)+(bc)+(cd)+(da)}{(2bc-3a+1)+(2cd-3b-1)+(2da-3c+1)+(2ab-3d-1)} \\
    &=\frac{ab+bc+cd+da}{2ab+2bc+2cd+2da-3a-3b-3c-3d+1-1+1-1} \\
    &=\frac{ab+bc+cd+da}{2(ab+bc+cd+da)-3(a+b+c+d)} \\
    &=\frac{ab+bc+cd+da}{2(ab+bc+cd+da)-3(0)} \\
    &=\frac{ab+bc+cd+da}{2(ab+bc+cd+da)} \\
    &=\frac{1}{2} \text{ } (\because ab+bc+cd+da\ne0)
\end{align*}
\\\\
\\\\
Conclusion: \textbf{The value of $k$ is $\frac{1}{2}$}.


\end{document}

